\subsection{A few hundred ligands recover the map; five-mode compression is exploratory}
\label{sec:cost}

Because the map depends on ligand support, it must be estimated on the
relevant chemical domain. We therefore asked how large a pilot sample drawn from the source library
must be docked before the estimate settles. We drew 100 random calibration
samples at each of seven sizes from 50 to 5,000 ligands and compared each
estimated residual map with the full-matrix map by edge-rank agreement
(Figure~\ref{fig:cost}a). Mean agreement for Docking-44 is 0.813 at 50 ligands,
0.942 at 200 and 0.976 at 500; DOCKSTRING-58 gives 0.766, 0.918 and 0.965 at the
same sizes. Most of the recovery occurs by 200 ligands. The next tenfold increase, from 500
to 5,000, adds only 0.022 on Docking-44 and 0.032 on DOCKSTRING-58.

Results are similar under a five-fold chemical-group holdout, run on all of Docking-44 and
on a fixed 15{,}000-ligand DOCKSTRING support, with calibration ligands drawn from
four folds and the reference map from the fifth: mean agreement at 200 ligands is
0.927 for Docking-44 and 0.916 for DOCKSTRING-58, rising to 0.959 and 0.958 at
500. In practical terms, a two-stage design at 200 ligands scores 11,600 of the
15,083,480 DOCKSTRING-58 cells, about 1,300-fold fewer; for Docking-44 it is 8,800
of 556,644, a 63-fold reduction. On the frozen 20-kinase retrieval endpoint of
Figure~\ref{fig:external}b, a 200-ligand sample gives mean ROC-AUC 0.701 and a
500-ligand sample gives 0.706, against 0.710 from all 260{,}060 DOCKSTRING ligands. These are
score-surface recovery diagnostics on the source library: a calibration sample recovers the
geometry of the library it came from.

An exploratory compression analysis asks whether the residual map is equally informative in every direction. On the 20-kinase
DOCKSTRING support, the smallest docking-only subspace holding at least half of the
residual trace has five modes, covering 0.517 of a trace spread over 19 positive
eigenvalues. The rule uses only the docking eigenvalues, with no experimental
values, and returns five in every one of five scaffold-disjoint folds;
across the ten fold pairs the mean squared canonical correlation between fold
subspaces runs from 0.986 to 0.996, indicating that the five-dimensional subspace
is stable across folds.

Truncating the residual correlation to those five modes and renormalising to unit
diagonal gives higher point estimates of rank agreement with the experimental target-pair map
on all four kinase panels, by 0.033 on PKIS2 up to 0.058 on KiRHub, and higher ROC-AUC point
estimates on all four, by 0.014 on PKIS1 up to 0.063 on KiRHub (Figure~\ref{fig:cost}b).
Average precision changes sign across panels: $+0.034$ on DAVIS and $+0.044$ on KiRHub, but
$-0.035$ on PKIS1 and $-0.027$ on PKIS2. The change is positive in 8 of 20 target deletions on
PKIS1 and 5 of 20 on PKIS2. The trace rule used no experimental value and the panels were
scored without retuning it, but PKIS1 was the discovery panel and the six-test validation
sensitivity was unresolved. Its interpretation across the four assay panels remains post hoc.

The participation ratio of the same residual map is 11.6, yet dropping the fourteen directions
outside this exploratory five-mode subspace gives higher point estimates of rank agreement and
ROC-AUC. Participation ratio counts directions, not their usefulness for a particular endpoint;
the spectrum and external-comparison analysis therefore answer different questions.
This is a within-scorer statement, and we do not extend it to comparisons between
scoring functions.

We release the map-construction and diagnostic stage of the audit as a command-line tool;
the transformation-matched null constructions are supplied as separate analysis scripts
rather than inside the CLI. \texttt{analysis/target\_map\_audit.py} takes any dense ligand-by-target score
matrix and writes the raw target-correlation map, the centred map, the spectrum of
each (participation-ratio dimension with its exact mean-squared-correlation
identity, PC1 variance fraction, entropy dimension, and components reaching 90\% of
trace), and a signed edge table. Optional arguments add the domain-quantile map
comparison of Figure~\ref{fig:residual}b and a calibration-recovery curve of the
kind in Figure~\ref{fig:cost}a, with each pilot map compared against its
row-disjoint complement. Each run records its preprocessing contract and SHA-256
hashes for every file written. The tool ships with
\texttt{analysis/test\_target\_map\_audit.py} and
\texttt{analysis/TARGET\_MAP\_AUDIT.md}.
